\chapter{Power tests and Moser iteration}\label{sec:moser-iteration}
We now prove the local-boundedness step used in elliptic regularity. Unless stated otherwise, $n\geq2$, $A$ is measurable and symmetric, and $\lambda I\leq A\leq\Lambda I$ with $0<\lambda\leq\Lambda$. A weak subsolution of $\operatorname{div}(A Du)=0$ satisfies $\int A Du\cdot D\phi\leq0$ for nonnegative compactly supported tests. There are no lower-order terms in this section. The source is \cite{SimonPDE}, Lecture 17, Theorem 1 and its preparatory lemmas, printed pp.~199--205.

\begin{theorem}[The Sobolev gain used in iteration]\label{thm:iteration-sobolev-gain}
For $n\geq2$, put $\chi=n/(n-2)$ if $n>2$, and $\chi=2$ if $n=2$. There is $S_n<\infty$ such that every $w\in H_0^1(B_2\subset\mathbb R^n)$ satisfies
\[
 \left(\int_{B_2}|w|^{2\chi}\right)^{1/\chi}
 \leq S_n\int_{B_2}|Dw|^2.
\]
See \cite{SimonPDE}, Lecture 17, Lemma 2, pp.~200--201, and Corollary 1, p.~202. The choice $n/(n-2)$ is not used in dimension two.
\end{theorem}
\begin{proof}
Corollary \ref{cor:scaled-sobolev-ball}, proved from the coordinate-projection Holder inequality and the first-power Sobolev estimate, gives the assertion with $R=2$. Its radius factor is fixed and can be included in $S_n$. In dimension two the proof uses the finite-power inequality of Theorem \ref{thm:two-dimensional-sobolev} with $q=4$.
\end{proof}

\begin{lemma}[Positive parts preserve the subsolution inequality]\label{lem:positive-part-subsolution}
If $v\in H^1_{\mathrm{loc}}(U)$ is a weak subsolution for the coefficients fixed above, then $(v-k)_+$ is a weak subsolution for every constant $k$. If $v$ is a weak solution, both $v_+$ and $(-v)_+$ are nonnegative weak subsolutions.
\end{lemma}
\begin{proof}
Choose smooth convex nondecreasing approximations $\beta_\varepsilon$ to the positive-part function, with $0\leq\beta'_\varepsilon\leq1$ and $\beta''_\varepsilon\geq0$. Test the subsolution inequality with $\eta\beta'_\varepsilon(v-k)$ for $0\leq\eta\in C_c^\infty(U)$. This is an admissible $H_0^1$ test; approximation by nonnegative smooth tests is justified because $A Dv$ is locally square integrable. The chain rule gives
\[
 \int \beta'_\varepsilon(v-k)A Dv\cdot D\eta
 \leq-\int\eta\beta''_\varepsilon(v-k)A Dv\cdot Dv\leq0.
\]
The Sobolev chain rule and the fact that $Dv=0$ almost everywhere on a level set give $\beta_\varepsilon(v-k)\to(v-k)_+$ in $H^1$ on compact subsets. Pass to the limit to get the required inequality. These elementary Sobolev composition and approximation facts are the functional-analytic inputs. For a solution, $-v$ is also a solution; apply the first assertion with $k=0$ to both signs. The constant shift is legitimate because the equation has no zero-order term.
\end{proof}

\begin{lemma}[Energy estimate for truncated powers]\label{lem:truncated-power-energy}
Let $u\geq0$ be an $H^1_{\mathrm{loc}}$ weak subsolution, let $p\geq2$, $L>0$, and set
\[
 u_L=\min\{u,L\},\qquad v_L=u\,u_L^{p/2-1}.
\]
For $p=2$ interpret $v_L=u$. For every nonnegative smooth compactly supported cutoff $\eta$,
\[
 \int|D(\eta v_L)|^2
 \leq C p^2(\Lambda/\lambda)^2\int v_L^2|D\eta|^2,
\]
where $C$ is an absolute constant independent of $p,L,u$.
\end{lemma}
\begin{proof}
For fixed $L$, the functions $s\mapsto s\min(s,L)^{p-2}$ and $s\mapsto s\min(s,L)^{p/2-1}$ are Lipschitz on $[0,\infty)$, with the $p=2$ case interpreted separately. Thus $\eta^2u u_L^{p-2}$ is an admissible nonnegative $H_0^1$ test, even though no boundedness of $u$ has been proved. The term obtained by differentiating $u_L^{p-2}$ is nonnegative by ellipticity. Dropping it and using Young's inequality gives
\[
 \int\eta^2 u_L^{p-2}|Du|^2
 \leq4(\Lambda/\lambda)^2\int u^2u_L^{p-2}|D\eta|^2.
\]
On $\{u<L\}$ the derivative of $v_L=u^{p/2}$ is $(p/2)u^{p/2-1}Du$, whereas on $\{u>L\}$ it is $L^{p/2-1}Du$. Consequently
\[
 |Dv_L|^2\leq(p^2/4)u_L^{p-2}|Du|^2
\]
almost everywhere. Combining this with $|D(\eta v_L)|^2\leq2\eta^2|Dv_L|^2+2v_L^2|D\eta|^2$ proves the estimate. Truncation gives an $H^1$ function with bounded value, not a claim of bounded gradient; this distinction matters in the source's power-test calculation, \cite{SimonPDE}, pp.~202--203.
\end{proof}

\begin{lemma}[One step of integrability improvement]\label{lem:moser-integrability-step}
Let $u\geq0$ be a weak subsolution on $B_2$, and suppose $u\in L^p(B_s)$ for $p\geq2$ and $1\leq r<s\leq2$. With $\chi$ as in Theorem \ref{thm:iteration-sobolev-gain},
\[
 \|u\|_{L^{p\chi}(B_r)}
 \leq\left(\frac{C p^2}{(s-r)^2}\right)^{1/p}\|u\|_{L^p(B_s)},
\]
where $C$ depends only on $n$ and $\Lambda/\lambda$.
\end{lemma}
\begin{proof}
Choose $\eta=1$ on $B_r$, supported in $B_s$, with $|D\eta|\leq C/(s-r)$. Apply the Sobolev gain to $\eta v_L$ and then Lemma \ref{lem:truncated-power-energy}. The right side is bounded by $Cp^2(s-r)^{-2}\int_{B_s}u^p$, because $v_L^2=u^2u_L^{p-2}\leq u^p$. As $L\to\infty$, Fatou's lemma gives
\[
 \left(\int_{B_r}u^{p\chi}\right)^{1/\chi}
 \leq\frac{Cp^2}{(s-r)^2}\int_{B_s}u^p.
\]
Taking the $p$th root proves the claim. In particular the higher integrability is a conclusion of the argument, not an assumption permitting an untruncated power test.
\end{proof}

\begin{theorem}[Moser local boundedness for nonnegative subsolutions]\label{thm:moser-subsolution-bound}
Let $n\geq2$, let $A$ be measurable and symmetric on $B_{2R}(a)$ with $\lambda I\leq A\leq\Lambda I$ for $0<\lambda\leq\Lambda$, and let $u\in H^1(B_{2R}(a))$ be a nonnegative weak subsolution of $\operatorname{div}(A Du)=0$. Then
\[
 \operatorname*{ess\,sup}_{B_R(a)}u
 \leq C(n,\Lambda/\lambda)\left(R^{-n}\int_{B_{2R}(a)}u^2\right)^{1/2}.
\]
This is the zero-lower-order, $p=2$ case of \cite{SimonPDE}, Lecture 17, Theorem 1, p.~200.
\end{theorem}
\begin{proof}
Translation and dilation reduce to $a=0,R=1$ without changing the ellipticity ratio. Let $r_j=1+2^{-j}$ and $p_j=2\chi^j$. Starting from $u\in L^2(B_2)$, Lemma \ref{lem:moser-integrability-step} successively gives
\[
 \|u\|_{L^{p_{j+1}}(B_{r_{j+1}})}
 \leq\bigl(Cp_j^2 2^{2j+2}\bigr)^{1/p_j}
 \|u\|_{L^{p_j}(B_{r_j})}.
\]
The logarithm of the product of these constants is bounded by a multiple of $\sum_j(1+j)\chi^{-j}$, which is finite. Thus $\|u\|_{L^{p_j}(B_1)}\leq C\|u\|_{L^2(B_2)}$ for every $j$. If $u>A$ on a subset of $B_1$ of positive measure $m$, then $\|u\|_{L^{p_j}(B_1)}\geq A m^{1/p_j}$. Letting $j\to\infty$ bounds every such level $A$ by $C\|u\|_2$, and hence bounds the essential supremum. This does not assume $u\in L^\infty$ in advance. Scaling the volume integral back gives the factor $R^{-n/2}$, including in dimension two.
\end{proof}

\begin{proposition}[Local bounds from every positive integrability exponent]\label{prop:moser-all-positive-exponents}
Under the hypotheses of Theorem \ref{thm:moser-subsolution-bound}, for every $p>0$,
\[
 \operatorname*{ess\,sup}_{B_R(a)}u
 \leq C(n,\Lambda/\lambda,p)\left(R^{-n}\int_{B_{2R}(a)}u^p\right)^{1/p}.
\]
The right side is allowed to be infinite. For $p<1$ no norm triangle inequality is needed. Compare \cite{SimonPDE}, Lecture 17, Theorem 1.
\end{proposition}
\begin{proof}
Work on $B_2$ by scaling. For $p\geq2$, Holder and Theorem \ref{thm:moser-subsolution-bound} suffice. For $0<p<2$, put $P=(\int_{B_2}u^p)^{1/p}$ and $M(r)=\operatorname*{ess\,sup}_{B_r}u$. Applying the half-ball estimate on small balls centered in $B_r$ gives, for $1\leq r<s\leq3/2$,
\[
 M(r)\leq C(s-r)^{-n/2}\|u\|_{L^2(B_s)}
 \leq C(s-r)^{-n/2}M(s)^{1-p/2}P^{p/2}.
\]
A finite covering of $B_r$ by balls of radius comparable to $s-r$ justifies the first inequality with the same constant. Young's inequality then yields, for any fixed $0<\theta<1$,
\[
 M(r)\leq\theta M(s)+C_{\theta,p}(s-r)^{-n/p}P.
\]
Choose $r_j=1+(1-2^{-j})/2$ and $\theta<2^{-n/p}$. Iteration gives a convergent geometric series for the error terms, since $r_{j+1}-r_j=2^{-j-2}$. The remainder is at most $\theta^j M(3/2)$ and tends to zero; $M(3/2)<\infty$ follows from the already proved local $L^2$ bound. Thus $M(1)\leq C_pP$. Rescaling proves the result. The argument uses interpolation of integrals and a scalar iteration, not Banach-space properties for $p<1$.
\end{proof}

\begin{proposition}[Weak maximum principle with a Sobolev boundary condition]\label{prop:weak-maximum-sobolev-boundary}
Let $\Omega\subset\mathbb R^n$ be bounded, $n\geq2$, and let $u\in H^1(\Omega)$ be a weak subsolution with the ellipticity assumptions above. If $(u-k)_+\in H_0^1(\Omega)$ for a constant $k$, then $u\leq k$ almost everywhere in $\Omega$. No regularity of the boundary of $\Omega$ is assumed; its boundary condition is expressed by this Sobolev membership.
\end{proposition}
\begin{proof}
Set $w=(u-k)_+$. It is an admissible nonnegative $H_0^1$ test, by approximation. The Sobolev chain rule gives $Dw=Du$ on $\{u>k\}$ and zero elsewhere, so
\[
 0\geq\int A Du\cdot Dw=\int A Dw\cdot Dw\geq\lambda\int|Dw|^2.
\]
Extend $w$ by zero to a ball containing $\Omega$. Its extension lies in $H_0^1$ of that ball. The scaled version of Theorem \ref{thm:iteration-sobolev-gain} now forces $w=0$, which is the assertion. Unlike a pointwise boundary statement, the assumed membership has meaning on any bounded open domain.
\end{proof}

\begin{corollary}[An integrable entire subsolution vanishes]\label{cor:integrable-subsolution-liouville}
Let $n\geq2$ and let $A$ be measurable and symmetric on all of $\mathbb R^n$, with fixed global bounds $\lambda I\leq A\leq\Lambda I$. If $u\geq0$ is a weak subsolution in $H^1_{\mathrm{loc}}(\mathbb R^n)$ and $\int_{\mathbb R^n}u^p<\infty$ for some $p>0$, then $u=0$ almost everywhere. A signed entire weak solution with finite $\int|u|^p$ also vanishes.
\end{corollary}
\begin{proof}
Proposition \ref{prop:moser-all-positive-exponents}, on balls of arbitrarily large radius, bounds $\operatorname*{ess\,sup}_{B_R}u$ by $C R^{-n/p}(\int_{\mathbb R^n}u^p)^{1/p}$, which tends to zero. Every fixed ball lies in these expanding balls, so $u$ vanishes on every ball. For a signed solution, its positive and negative parts are subsolutions by Lemma \ref{lem:positive-part-subsolution}, and each has finite $p$ integral. Apply the first assertion to both. The global ellipticity bounds and the integrability hypothesis are essential parts of this conclusion.
\end{proof}
